% part: intuitionistic-logic
% chapter: propositions-as-types
% section: sequent-natural-deduction

\documentclass[../../../include/open-logic-section]{subfiles}

\begin{document}

\olfileid[tr]{int}{pty}{snd}

\olsection{Sekant Biçiminde Doğal Tümdengelim}

!!{undischarged} varsayımları $\Gamma$, sonucu $!A$ olan bir doğal
tümdengelim !!{derivation}{}i varsa $\Gamma \Sequent !A$; $\Gamma$ boşsa
$\Sequent !A$ yazalım.

$\Gamma \cup \{!A_1, \dots, !A_n\}$ yerine $\Gamma, !A_1, \dots, !A_n$;
$\Gamma \cup \Delta$ yerine $\Gamma, \Delta$ yazıyoruz.

$\Gamma \Sequent !A \land !B$ olduğunda, yani !!{undischarged}
varsayımları $\Gamma$ ve son-!!{formula}{}ü $!A \land !B$ olan bir
!!a{derivation} bulunduğunda, en alta $\Elim{\land}$ uygulayarak aynı
!!{undischarged} varsayımlara ve sonuç olarak $!A$'ya sahip bir
!!a{derivation} elde edebiliriz. Başka bir deyişle $\Gamma \Sequent !A
\land !B$ ise $\Gamma \Sequent !A$ olur.
\begin{prooftree}
  \Axiom$\Gamma \fCenter !A \land !B$
  \RightLabel{$\Elim{\land}$}
  \UnaryInf$\Gamma \fCenter !A$
  \DisplayProof\qquad\bottomAlignProof
  \Axiom$\Gamma \fCenter !A \land !B$
  \RightLabel{$\Elim{\land}$}
  \UnaryInf$\Gamma \fCenter !B$
\end{prooftree}
$\Elim{\land}$ etiketi, doğal tümdengelimdeki aynı adlı kuralla ilişkiye
işaret eder.

Benzer biçimde $\Gamma, !A \Sequent !B$ olduğunu, yani !!{undischarged}
varsayımları $\Gamma, !A$ ve son-!!{formula}{}ü~$!B$ olan bir
!!a{derivation} bulunduğunu varsayın. $\Intro{\lif}$ kuralını uygularsak
!!{undischarged} varsayımları $\Gamma$ ve son-!!{formula}{}ü $!A \lif !B$
olan bir !!a{derivation} elde ederiz; yani $\Gamma \Sequent !A \lif !B$.
Bunun varsayımların !!{discharge} işlemini nasıl daha açık kıldığına dikkat
edin.
\begin{prooftree}
  \Axiom $\Gamma, !A \fCenter !B$
  \RightLabel{$\Intro{\lif}$}
  \UnaryInf $\Gamma \fCenter !A \lif !B$
\end{prooftree}

Diğer kurallardan da aynı biçimde sonuçlar çıkarabiliriz; bunlar açıkça
şöyle yazılır:
\begin{gather*}
  \Axiom $\Gamma \fCenter !A$
  \Axiom $\Delta \fCenter !B$
  \RightLabel{$\Intro{\land}$}
  \BinaryInf $\Gamma,\Delta \fCenter !A \land !B$
  \DisplayProof\\
  \Axiom $\Gamma \fCenter !A \land !B$
  \RightLabel{$\Elim{\land}_1$}
  \UnaryInf $\Gamma \fCenter !A$
  \DisplayProof
  \qquad
  \Axiom $\Gamma \fCenter !A \land !B$
  \RightLabel{$\Elim{\land}_2$}
  \UnaryInf $\Gamma \fCenter !B$
  \DisplayProof
  \\
  \Axiom $\Gamma \fCenter !A$
  \RightLabel{$\Intro{\lor}_1$}
  \UnaryInf $\Gamma \fCenter !A \lor !B$
  \DisplayProof\qquad
  \Axiom $\Gamma \fCenter !B$
  \RightLabel{$\Intro{\lor}_2$}
  \UnaryInf $\Gamma \fCenter !A \lor !B$
  \DisplayProof\\
  \Axiom $\Gamma \fCenter !A \lor !B$
  \Axiom $\Delta, !A \fCenter !C$
  \Axiom $\Delta', !B \fCenter !C$
  \RightLabel{$\Elim{\lor}$}
  \TrinaryInf $\Gamma, \Delta, \Delta' \fCenter !C$
  \DisplayProof
  \\
  \Axiom $\Gamma, !A \fCenter !B$
  \RightLabel{$\Intro{\lif}$}
  \UnaryInf $\Gamma \fCenter !A \lif !B$
  \DisplayProof
  \qquad
  \Axiom $\Delta \fCenter !A \lif !B$
  \Axiom $\Gamma \fCenter !A$
  \RightLabel{$\Elim{\lif}$}
  \BinaryInf $\Gamma, \Delta \fCenter !B$
  \DisplayProof
  \\
  \Axiom $\Gamma \fCenter \lfalse$
  \RightLabel{$\FalseInt$}
  \UnaryInf $\Gamma \fCenter !A$
  \DisplayProof
\end{gather*}

Her varsayım tek başına $!A$'nın~$!A$'dan bir !!a{derivation}{}idir; yani her
zaman $!A \Sequent !A$ olur.

\begin{prooftree}
  \AxiomC{$ $}
  \UnaryInfC{$!A \fCenter !A$}
\end{prooftree}

Bu kurallar birlikte, hangi doğal tümdengelim !!{derivation}s{}inin var
olduğunu belirleyen bir hesap olarak alınabilir. Ayrıca her adımda yalnız
!!{derive} eylemiyle elde edilen !!{formula} değil, onun hangi
!!{undischarged} varsayımlardan !!{derive} eylemiyle elde edildiği de
kaydedildiğinden doğal
tümdengelimin gösterimsel bir
çeşitlemesi olarak da görülebilirler.

\begin{prooftree}
  \Axiom$!A \fCenter !A$
  \RightLabel{$\Intro{\lor}_1$}
  \UnaryInf $!A \fCenter !A \lor (!A \lif \lfalse)$
  \Axiom $!B \fCenter !B$
  \RightLabel{$\Elim{\lif}$}
  \BinaryInf$!A, !B \fCenter \lfalse$
  \RightLabel{$\Intro{\lif}$}
  \UnaryInf$!B \fCenter !A \lif \lfalse$
  \RightLabel{$\Intro{\lor}_2$}
  \UnaryInf$!B \fCenter !A \lor (!A \lif \lfalse)$
  \Axiom$!B \fCenter !B$
  \RightLabel{$\Elim{\lif}$}
  \BinaryInf$!B \fCenter \lfalse$
  \RightLabel{$\Intro{\lif}$}
  \UnaryInf$\fCenter !B \lif \lfalse$
\end{prooftree}
burada $!B$, $(!A \lor (!A \lif \lfalse))\lif \lfalse$'ın kısaltmasıdır.

\end{document}
